% Author contribution and AI-use disclosure. Placement: after the conclusions, before the
% bibliography.
%
% Rewritten 2026-09-01 against the session archive. Two deliberate departures from the
% earlier version, which was instantiated from the hub's writing-style-and-ai-use.md §2:
%
%   1. The machine-checked/axiom-ledger paragraph is GONE. §1.1 (sec:machine-checked) owns
%      the trust boundary and states it better; repeating it here padded the statement and
%      buried the disclosure proper.
%   2. The guide's framing --- AI "assisting the author to re-present and verify his own
%      prior results" --- understates what happened, and the author's own account of the work
%      contradicts it. The hemigroup formulation and much of the mathematics postdate
%      2026-08-05. Claiming less AI involvement than there was is the one failure mode a
%      disclosure cannot afford.
%
% The "By the numbers" figures come from the session archive (ai-archive / chronicler). They
% have been revised four times as missing material was recovered — always upward, never down:
% 265 h -> 302 -> 305 -> 330. The 2026-09-05 revision followed the cloud-staging mechanism
% landing that session's full transcript with its sidecars. v1.1.0 printed 520 h (computed
% 2026-09-27, window April to 27 September); re-derived 2026-10-02 over the SAME window, after a
% full `chronicler run`: `chronicler stats --research --since 2026-04-01 --until 2026-09-27`
% gives 559.1 h, 105 884 typed words, 1 873 611 assistant words (17.7x), 86 302 directed words;
% the Lean week is `--since 2026-08-09 --until 2026-08-15` (205 human turns, 9 710 typed words),
% and its lines are git's (`--numstat` on Formalization/*.lean over those days: 24 698 added;
% 21 772 lines standing at the end of 15 August). The window is the work up to this article's
% v1.1.0 (the author, 2026-10-02), so later work on other papers is not counted. Do NOT restate them without
% re-deriving; `chronicler stats` is the source, and the computation date now appears in the
% prose because these numbers have a shelf life.
%
% "How the work went" deliberately does NOT say that months of reading preceded any writing —
% an earlier draft said that and it was false. There was a great deal of writing first (the
% research wiki, and the larger project's article and Lean development); what was missing was
% the way to weaken the axioms, and the hemigroup formulation is what unstuck it. Do not
% compress this back into a run-up: the stuck-then-unstuck shape is the true one.
%
% Timeline note: git dates the draft to 2026-08-07, which is the bootstrap commit, not the
% writing. The archive puts the draft's own prose in the 08-05 conversation --- "self-
% decomposable" at 21:05 that evening, "a measurement of a measurement is a measurement" at
% 07:06 the next morning. The prose here follows the archive, not the commit dates.

\section*{Author contribution and use of AI}

This article was written in close collaboration with generative AI, and that collaboration
went well beyond editing assistance. The disclosure below is specific because a general one
would misrepresent it.

\paragraph{What the author brought.} The research programme, the questions it asks, and the
point of departure are the author's. The semigroup theory of \sscite{fagerstrom2005temporal}
is his own earlier work, and the decision to ask what survives of it when the semigroup axiom
is weakened to a hemigroup is his. Every structural choice about what this article is, what
it claims, and what it declines to claim was made by him.

\paragraph{How the work went.} This material was to have formed part of a larger project, and
for months before this article existed the author was reading, building tools, accumulating a
research wiki, and writing that other article and its formalization --- held up on a single
point. The semigroup axiom plainly had to be weakened, and the right way to weaken it had not
been found. The hemigroup formulation was that way, and it arrived on the evening of 5 August
2026. The substance of the working draft --- the axioms, the characterization, and the shape
of the argument that carries them --- was written in that same extended conversation, overnight
and into the following day. The article repository was opened two days later; the Lean~4
development was produced over the seven days after that. The remaining three weeks went on what
the author takes to be the substance of authorship: understanding the mathematics well enough
to stand behind it, replacing serviceable argument with motivated exposition, supplying what
was missing, and scoping claims that had outrun their proofs.

\paragraph{What the AI did.} The prose and the code were written by AI agents throughout, in
dialogue with the author, who worked almost entirely from the conversation rather than from
the editor. This includes the Lean~4 formalization essentially in full: the author's
involvement there was slight, and none of the proof terms are his. It also includes
mathematical content. A number of the results, the connections to the probabilistic
literature, and several of the lines of argument were not in the author's plan and were
learnt from the agents in the course of the work. They belong to the article, not merely to
its presentation.

\paragraph{AI was also the first reviewer.} Adversarial review was a working method
throughout, not only a final gate: drafts were read blind, section by section, by reviewer
agents that had no part in the writing, and the content revision was worked from their
findings. Before the final revision the manuscript was then read by AI systems other than
those that had written it: one review in journal style, returning a verdict of major
revision, and one on presentation. Their findings drove the last pass --- several headline
claims were scoped back to what is actually proved, one proof was repaired, and one
objection was answered with a new result. Both of these reviews are archived verbatim in
the repository, under \texttt{notes/reviews/}, beside the plan that records the disposition
of each finding. The author assessed each finding independently and accepted, narrowed, or
declined it; no change was made because a reviewer asked for it alone. The author has read
and revised the article throughout; beyond him, no human reader has yet reviewed it.

\paragraph{By the numbers.} The sessions were archived as the work went on, so the division of
labour can be stated rather than estimated. Between April and September 2026 the work ran to
some 560 hours of recorded active session time. The author contributed about 106\,000 words of
direction and the agents about 1.87 million words of response, a ratio of roughly eighteen to
one. The seven days that produced the Lean development wrote 24\,700 lines of it, and took 205
turns and 9\,710 words from the author --- under four tenths of a word per line of proof.

These totals cover the work up to 27 September 2026, when version 1.1.0 was released; they are
lower bounds, and they were computed on 2 October 2026. Both qualifications are load-bearing. A handful of sessions leave no transcript; some subagent traces were lost to
context compaction before they could be archived; and reading on paper, thinking away from a
keyboard, and work done under a separate professional obligation whose records cannot be
exported leave no trace at all. The figures have also only ever moved in one direction as
missing material was recovered --- by roughly a fifth over the week the first version of this
statement was written, and by about seven per cent more when they were computed again for this
version over the same period --- so they understate rather than overstate the agents' share.

Two conventions, because neither is self-evident. The line count is of lines of Lean \emph{written}
in those seven days, not of lines standing at the end of them: the development stood at 21\,800
lines on 15 August 2026, having been edited as well as extended. And the author's word count
excludes the task briefs of unattended sessions, which this programme has dispatched from a
work queue since 21 September 2026: those briefs are written by an agent, not typed by the
author, and counting them as direction would have inflated it by some four fifths and made
the ratio read as ten to one.

\paragraph{What this article would otherwise have been.} The author would not, realistically,
have written this article unaided --- certainly not under the constraints of research pursued
as a spare-time activity. Written alone it would have carried substantially less mathematical
detail, and would have leaned on results imported from the probability literature at many of
the points where it now proves them.

\paragraph{Responsibility.} The author directed the work throughout, has satisfied himself of
the mathematical content, and has checked every citation. Responsibility for this article,
including for its errors, is his alone. What is machine-checked, and the axioms on which the
verified development rests, are set out in \S\ref{sec:machine-checked}; the validity of those
results depends on the Lean kernel and not on how, or by whom, the proof terms were produced.
